\bibitem{s1ref1} Y. Yan et al. Magma evolution and mineralization of the Baixintan magmatic Ni--Cu sulfide deposit in Eastern Tianshan, Northwestern China. \textit{Int. J. Earth Sci.}, 111:2823--2843, 2022. \doi{10.1007/s00531-022-02169-2}.

\bibitem{s1ref2} R. G. Cawthorn. The Platinum Group Element Deposits of the Bushveld Complex in South Africa. \textit{Platinum Met. Rev.}, 54:205--215, 2010. \doi{10.1595/147106710x520222}.

\bibitem{s1ref3} R. Dahl. Geology of the Two Duck Lake Intrusion and the Marathon Cu-PGE Deposit, Coldwell Complex, Northern Ontario. \textit{Explor. Min. Geol.}, 10:51--65, 2001. \doi{10.2113/10.1-2.51}.

\bibitem{s1ref4} A. A. Ariskin et al. Cu--Ni--PGE fertility of the Yoko-Dovyren layered massif (northern Transbaikalia, Russia): thermodynamic modeling of sulfide compositions in low mineralized dunite based on quantitative sulfide mineralogy. \textit{Miner. Deposita}, 51:993--1011, 2016. \doi{10.1007/s00126-016-0666-8}.

\bibitem{s1ref5} E. M. Ripley. Ni-Cu-PGE Mineralization in the Partridge River, South Kawishiwi, and Eagle Intrusions: A Review of Contrasting Styles of Sulfide-Rich Occurrences in the Midcontinent Rift System. \textit{Econ. Geol.}, 109:309--324, 2013. \doi{10.2113/econgeo.109.2.309}.

\bibitem{s1ref6} M. Sciortino, J. E. Mungall, and J. Muinonen. Generation of High-Ni Sulfide and Alloy Phases During Serpentinization of Dunite in the Dumont Sill, Quebec. \textit{Econ. Geol.}, 110:733--761, 2015. \doi{10.2113/econgeo.110.3.733}.

\bibitem{s1ref7} D. J. McKevitt, C. M. Lesher, and M. G. Houlé. Volcanology, geochemistry and petrogenesis of the Expo-Raglan magmatic system in the eastern Cape Smith Belt, Nunavik, northern Quebec. \textit{Natural Resources Canada/CMSS/Information Management}, 2019. \doi{10.4095/313674}.

\bibitem{s1ref8} T. Oberthür. Platinum-Group Element Mineralization of the Main Sulfide Zone, Great Dyke, Zimbabwe. \textit{Magmatic Ni-Cu and PGE Deposits: Geology, Geochemistry, and Genesis}, 2011. \doi{10.5382/rev.17.12}.

\bibitem{s1ref9} K. Kuster, C. M. Lesher, and M. G. Houlé. Geology and geochemistry of mafic and ultramafic bodies in the Shebandowan mine area, Wawa-Abitibi terrane: implications for Ni-Cu-(PGE) and Cr-(PGE) mineralization, Ontario and Quebec. \textit{Natural Resources Canada/CMSS/Information Management}, 2022. \doi{10.4095/329394}.

\bibitem{s1ref10} D. Tang, K. Qin, B. Su, P. A. Sakyi, Y. Mao, and S. Xue. Petrogenesis and mineralization of the Hulu Ni-Cu sulphide deposit in Xinjiang, NW China: constraints from Sr-Nd isotopic and PGE compositions. \textit{Int. Geol. Rev.}, 56:711--733, 2014. \doi{10.1080/00206814.2014.893413}.

\bibitem{s1ref11} J. Lehmann, N. Arndt, B. Windley, M. F. Zhou, C. Y. Wang, and C. Harris. Field Relationships and Geochemical Constraints on the Emplacement of the Jinchuan Intrusion and its Ni-Cu-PGE Sulfide Deposit, Gansu, China. \textit{Econ. Geol.}, 102:75--94, 2007. \doi{10.2113/gsecongeo.102.1.75}.

\bibitem{s1ref12} E. J. Hanski, Z. Y. Luo, H. Oduro, and R. J. Walker. The Pechenga Ni-Cu Sulfide Deposits, Northwestern Russia: A Review with New Constraints from the Feeder Dikes. \textit{Magmatic Ni-Cu and PGE Deposits: Geology, Geochemistry, and Genesis}, 2011. \doi{10.5382/rev.17.05}.

\bibitem{s1ref13} S. J. Barnes, V. A. Melezhik, and S. V. Sokolov. The composition and mode of formation of the Pechenga nickel deposits, Kola Peninsula, northwestern Russia. \textit{Can. Mineral.}, 39:447--471, 2001. \doi{10.2113/gscanmin.39.2.447}.

\bibitem{s1ref14} C. L. Seabrook, H. M. Prichard, and P. C. Fisher. PLATINUM-GROUP MINERALS IN THE RAGLAN Ni Cu (PGE) Sulfide deposit, Cape Smith, Quebec, Canada. \textit{Can. Mineral.}, 42:485--497, 2004. \doi{10.2113/gscanmin.42.2.485}.

\bibitem{s1ref15} F. Santaguida, K. Luolavirta, M. Lappalainen, J. Ylinen, T. Voipio, and S. Jones. The Kevitsa Ni-Cu-PGE Deposit in the Central Lapland Greenstone Belt in Finland. \textit{Mineral Deposits of Finland}, 195--210, 2015. \doi{10.1016/b978-0-12-410438-9.00008-x}.

\bibitem{s1ref16} M. Iljina, W. Maier, and T. Karinen. PGE-(Cu-Ni) Deposits of the Tornio-Näränkävaara Belt of Intrusions (Portimo, Penikat, and Koillismaa). \textit{Mineral Deposits of Finland}, 133--164, 2015. \doi{10.1016/b978-0-12-410438-9.00005-4}.

\bibitem{s1ref17} E. Mansur et al. Formation of the Ertelien and Langedalen magmatic Ni--Cu sulfide deposits in Norway: investigating the evolution of platinum-group-element-depleted systems at convergent margins. \textit{European Journal of Mineralogy}, 37:841--869, 2025. \doi{10.5194/ejm-37-841-2025}.

\bibitem{s1ref18} A. Beinlich, A. von Heydebrand, R. Klemd, L. Martin, and J. Hicks. Desulphurisation, chromite alteration, and bulk rock PGE redistribution in massive chromitite due to hydrothermal overprint of the Panton Intrusion, east Kimberley, Western Australia. \textit{Ore Geol. Rev.}, 118:103288, 2020. \doi{10.1016/j.oregeorev.2019.103288}.

\bibitem{s1ref19} N. Norbu, J. Li, Z. Zhang, and W. Zhang. The Late Silurian Langmuri Cu-Ni Deposit in East Kunlun Orogenic Belt: An Example of Magmatic Sulfide Deposit in a Post-Collisional Extensional Setting. \textit{Isl. Arc}, 34, 2025. \doi{10.1111/iar.70013}.

\bibitem{s1ref20} Y. Tao, C. Li, X. Y. Song, and E. M. Ripley. Mineralogical, petrological, and geochemical studies of the Limahe mafic--ultramatic intrusion and associated Ni--Cu sulfide ores, SW China. \textit{Miner. Deposita}, 43:849--872, 2008. \doi{10.1007/s00126-008-0207-1}.

\bibitem{s1ref21} E. T. Mansur and C. F. Ferreira Filho. Magmatic structure and geochemistry of the Luanga Mafic--Ultramafic Complex: Further constraints for the PGE-mineralized magmatism in Carajás, Brazil. \textit{Lithos}, 266-267:28--43, 2016. \doi{10.1016/j.lithos.2016.09.036}.

\bibitem{s1ref22} V. Bécu and M. G. Houlé. Chromite occurrences of the E-Ext stripped outcrop, Mayville intrusion, Bird River greenstone belt, southeastern Manitoba. \textit{Natural Resources Canada/CMSS/Information Management}, 2015. \doi{10.4095/296135}.

\bibitem{s1ref23} E. V. Sharkov and A. V. Chistyakov. Geological and petrological aspects of Ni-Cu-PGE mineralization in the early Paleoproterozoic Monchegorsk layered mafic-ultramafic complex, Kola Peninsula. \textit{Geol. Ore Deposits}, 56:147--168, 2014. \doi{10.1134/s1075701514030040}.

\bibitem{s1ref24} N. T. Arndt, G. K. Czamanske, R. J. Walker, C. Chauvel, and V. A. Fedorenko. Geochemistry and Origin of the Intrusive Hosts of the Noril'sk-Talnakh Cu-Ni-PGE Sulfide Deposits. \textit{Econ. Geol.}, 98:495--515, 2003. \doi{10.2113/gsecongeo.98.3.495}.

\bibitem{s1ref25} S. J. Barnes, V. Taranovic, J. M. Miller, G. Boyce, and S. Beresford. Sulfide Emplacement and Migration in the Nova-Bollinger Ni-Cu-Co Deposit, Albany-Fraser Orogen, Western Australia. \textit{Econ. Geol.}, 115:1749--1776, 2020. \doi{10.5382/econgeo.4758}.

\bibitem{s1ref26} S. J. Barnes, D. R. Mole, R. Hornsey, and L. E. Schoneveld. Nickel-Copper Sulfide Mineralization in the Ntaka Hill Ultramafic Complex, Nachingwea Region, Tanzania. \textit{Econ. Geol.}, 114:1135--1158, 2019. \doi{10.5382/econgeo.4677}.

\bibitem{s1ref27} D. A. Gabov. PGE and Au minerals from the low-sulfide ore of the Pana Tundra Pluton, Kola Peninsula. \textit{Geol. Ore Deposits}, 52:624--629, 2010. \doi{10.1134/s1075701510070147}.

\bibitem{s1ref28} C. Li, E. M. Ripley, and Y. Tao. Chapter 11 Magmatic Ni-Cu and Pt-Pd Sulfide Deposits in China. \textit{Mineral Deposits of China}, 483--508, 2019. \doi{10.5382/sp.22.11}.

\bibitem{s1ref29} Y. Liu, W. Li, X. Lü, Y. Liu, B. Ruan, and X. Liu. Sulfide saturation mechanism of the Poyi magmatic Cu-Ni sulfide deposit in Beishan, Xinjiang, Northwest China. \textit{Ore Geol. Rev.}, 91:419--431, 2017. \doi{10.1016/j.oregeorev.2017.09.013}.

\bibitem{s1ref30} S. G. Todd, D. W. Keith, L. W. Le Roy, D. J. Schissel, E. L. Mann, and T. N. Irvine. The J-M platinum-palladium reef of the Stillwater Complex, Montana; I, Stratigraphy and petrology. \textit{Econ. Geol.}, 77:1454--1480, 1982. \doi{10.2113/gsecongeo.77.6.1454}.

\bibitem{s1ref31} A. Y. Barkov, J. G. Laflamme, L. J. Cabri, and R. F. Martin. Platinum-group minerals from the Wellgreen Ni-Cu-PGE deposit, Yukon, Canada. \textit{Can. Mineral.}, 40:651--669, 2002. \doi{10.2113/gscanmin.40.2.651}.

\bibitem{s1ref32} H. L. He et al. Genesis of the Xiarihamu Magmatic Ni-Co Sulfide Deposit in the East Kunlun Orogen, Northern Tibetan Plateau: In Situ Oxygen Isotope and Geochemical Perspectives. \textit{Econ. Geol.}, 117:1827--1844, 2022. \doi{10.5382/econgeo.4949}.

\bibitem{s1ref33} E. Farahat and H. Helmy. Abu Hamamid Neoproterozoic Alaskan-type complex, south Eastern Desert, Egypt. \textit{J. Afr. Earth Sci.}, 45:187--197, 2006. \doi{10.1016/j.jafrearsci.2006.02.003}.

\bibitem{s1ref34} D. A. Holwell, R. R. Keays, E. A. Firth, and J. Findlay. Geochemistry and mineralogy of platinum group element mineralization in the River Valley Intrusion, Ontario, Canada: a model for early-stage sulfur saturation and multistage emplacement and the implications for "contact-type" Ni-Cu-PGE sulfide mineralization. \textit{Econ. Geol.}, 109:689--712, 2014. \doi{10.2113/econgeo.109.3.689}.

\bibitem{s1ref35} E. Oliveira and J. Tarney. Genesis of the Precambrian copper-rich Caraiba hypersthenite-norite complex, Brazil. \textit{Miner. Deposita}, 30, 1995. \doi{10.1007/bf00202279}.

\bibitem{s1ref36} I. R. Rakhimov and V. N. Puchkov. Genesis of Sulfide and Platinum-Metal Mineralization of the Khudolaz Complex (Southern Urals). \textit{Dokl. Earth Sci.}, 498:398--402, 2021. \doi{10.1134/s1028334x21050135}.

\bibitem{s1ref37} N. F. Pchelintseva and E. V. Koptev-Dvornikov. On the genesis of platiniferous low-sulfide rocks in the Kivakka layered intrusion, Northern Karelia. \textit{Geochem. Int.}, 46:627--633, 2008. \doi{10.1134/s0016702908060098}.

\bibitem{s1ref38} M. Tarkian and T. Mutanen. Loveringite from the koitelainen layered intrusion, Northern Finland. \textit{Mineral. Petrol.}, 37:37--50, 1987. \doi{10.1007/bf01163157}.

\bibitem{s1ref39} A. Vymazalová, F. Laufek, T. Grokhovskaya, and C. Stanley. Viteite, Pd5InAs, a new mineral from the Monchetundra layered intrusion, Kola Peninsula, Russia. \textit{Can. Mineral.}, 58:395--402, 2020. \doi{10.3749/canim.1900080}.

\bibitem{s1ref40} L. Hulbert. Geology of the Muskox Intrusion and associated Ni+Cu occurrences. \textit{Natural Resources Canada/CMSS/Information Management}, 2005. \doi{10.4095/220361}.

\bibitem{s1ref41} V. Järvinen, T. Halkoaho, J. Konnunaho, J. S. Heinonen, and O. T. Rämö. Parental magma, magmatic stratigraphy, and reef-type PGE enrichment of the 2.44-Ga mafic-ultramafic Näränkävaara layered intrusion, Northern Finland. \textit{Miner. Deposita}, 55:1535--1560, 2019. \doi{10.1007/s00126-019-00934-z}.

\bibitem{s1ref42} N. Tolstykh, A. Kozlov, and Y. Telegin. Platinum mineralization of the Svetly Bor and Nizhny Tagil intrusions, Ural Platinum Belt. \textit{Ore Geol. Rev.}, 67:234--243, 2015. \doi{10.1016/j.oregeorev.2014.12.005}.

\bibitem{s1ref43} T. A. A. Halkoaho, T. T. Alapieti, J. J. Lahtinen, and J. M. Lerssi. The Ala-Penikka PGE reefs in the Penikat layered intrusion, Northern Finland. \textit{Mineral. Petrol.}, 42:23--38, 1990. \doi{10.1007/bf01162682}.

\bibitem{s1ref44} J. C. O. Andersen, H. Rasmussen, T. F. D. Nielsen, and J. G. Ronsbo. The Triple Group and the Platinova gold and palladium reefs in the Skaergaard Intrusion; stratigraphic and petrographic relations. \textit{Econ. Geol.}, 93:488--509, 1998. \doi{10.2113/gsecongeo.93.4.488}.

\bibitem{s1ref45} S. F. Sluzhenikin and N. A. Krivolutskaya. Pyasino--Vologochan intrusion: Geological structure and platinum--copper--nickel ores (Norilsk Region). \textit{Geol. Ore Deposits}, 57:381--401, 2015. \doi{10.1134/s1075701515050050}.

\bibitem{s1ref46} F. Villares et al. Petrogenesis of the Tampanchi Ultramafic--Mafic Complex (Ecuador): Geodynamic implications for the northwestern margin of South America during the late Cretaceous. \textit{Gondwana Research}, 105:514--534, 2022. \doi{10.1016/j.gr.2021.10.005}.

\bibitem{s1ref47} Y. Deng, F. Yuan, T. Zhou, C. Xu, D. Zhang, and X. Guo. Geochemical characteristics and tectonic setting of the Tuerkubantao mafic-ultramafic intrusion in West Junggar, Xinjiang, China. \textit{Geosci. Front.}, 6:141--152, 2015. \doi{10.1016/j.gsf.2013.10.003}.

\bibitem{s1ref48} H. Zhong et al. Trace-element and Sr--Nd isotopic geochemistry of the PGE-bearing Xinjie layered intrusion in SW China. \textit{Chem. Geol.}, 203:237--252, 2004. \doi{10.1016/j.chemgeo.2003.10.008}.

